\section{BCA-Style Accounting Logic}

The accounting exercise follows the classical BCA sequence: infer wedges from
data or simulated allocations, then run counterfactuals by allowing one wedge at
a time to move while holding the others fixed.

\subsection{Step 1: Infer Canonical Wedges}

Given series for \(\{Y_t,C_t,X_t,K_t,L_t\}\), infer the wedges as follows:
\begin{align}
    \widehat{A}_t
    &=
    \frac{Y_t}{F(K_t,L_t)},
    \label{eq:infer_efficiency}
    \\
    \widehat{G}_t
    &=
    Y_t-C_t-X_t,
    \label{eq:infer_government}
    \\
    1-\widehat{\tau}_{\ell,t}
    &=
    \frac{-U_L(C_t,L_t)}
    {U_C(C_t,L_t)\widehat{A}_tF_L(K_t,L_t)},
    \label{eq:infer_labor}
    \\
    \widehat{\xi}^x_t
    &=
    \E_t m^k_{t+1}.
    \label{eq:infer_investment_log}
\end{align}
Equation \eqref{eq:infer_investment_log} writes the investment wedge as a log
Euler residual. Under the second-order no-friction benchmark, this residual is
\[
    \widehat{\xi}^x_t
    =
    -\frac{1}{2}\Var_t(m^k_{t+1}).
\]
In empirical BCA, expectations and the stochastic process for wedges are usually
estimated jointly. This note isolates the algebraic interpretation.

\subsection{Step 2: Test the TFP-Uncertainty Restriction}

The single-source TFP uncertainty hypothesis implies
\begin{equation}
    \widehat{\xi}^x_t
    =
    \lambda_x q_t + u^x_t,
    \label{eq:test_investment_loading}
\end{equation}
where \(q_t\) is the centered variance state of the efficiency wedge. A strong
relationship between \(\widehat{\xi}^x_t\) and \(q_t\) would support the idea
that second-moment movements in the efficiency wedge show up as investment
wedges. Large systematic movements in labor or resource wedges would instead
suggest that the model is missing other mechanisms or that the data require
additional primitive shocks.

\subsection{Step 3: Counterfactual Wedge Experiments}

The counterfactual exercises are the standard BCA exercises:
\begin{enumerate}[label=(\alph*)]
    \item Efficiency-only: allow \(\widehat{A}_t\) and its volatility state
    \(q_t\) to move; hold other wedges fixed.
    \item Investment-wedge channel: allow the investment wedge implied by
    Equation \eqref{eq:investment_uncertainty_wedge} to move; hold labor and
    resource wedges fixed.
    \item Labor-wedge channel: allow \(\widehat{\tau}_{\ell,t}\) to move; hold
    efficiency, investment, and resource wedges fixed.
    \item Resource-wedge channel: allow \(\widehat{G}_t\) to move; hold other
    wedges fixed.
\end{enumerate}
The contribution of each wedge is evaluated by how closely the counterfactual
paths reproduce the observed paths of output, consumption, investment, and
labor.

\subsection{Connection to the Later SOE Nonlinear BCA}

The later small-open-economy framework will expand the resource wedge into
trade and external-finance objects, and may add a UIP or country-premium wedge.
The present closed-economy note establishes the core logic first:
\[
    \begin{aligned}
    \text{stochastic volatility in the efficiency wedge}
    &\rightarrow
    \text{Euler-equation risk correction}
    \\
    &\rightarrow
    \text{investment wedge representation}.
    \end{aligned}
\]
